% Provenance: copied from the design workspace CambridgePQ/recursion.fast, `report/sections/14-issues-roadmap.tex`, on 2026-10-19. Read-only copy; the source is authoritative.
\section{Issue register and roadmap}\label{sec:issues}

This section is the working list for the organising team. It collects every gap, objection and
fix proposal raised by the design documents (D1--D4), the adversarial reviews (A1--A3 of the design,
A4 of the compiled v0.2 contract, A5 of this report's cryptography, A6 of the compiled v0.3 contract
and its harness, A7 and A8 of this report's formal methods and editing), the eight research notes (R1--R8), the baseline (\cref{sec:baseline}), the critique
of the brief and the open items of DESIGN, deduplicates them, and records for each one what is to be
done, by whom and when. Most findings of the reviews are already fixed in the specification in force
(DESIGN v0.3 with the round-3 amendments of §15.4) and in the compiled contract; they are listed
compactly in \cref{subsec:iss-resolved} so that the reader can trace each attack to the rule that
answers it. The open issues follow in full in \cref{subsec:iss-open}, and the phased plan that
schedules them is in \cref{subsec:iss-roadmap}.

% local helpers for the issue boxes (prefixed to avoid clashes with other sections)
\newcommand{\IRprio}[1]{\textbf{\textcolor{red!55!black}{#1}}}
\newcommand{\IRprob}{\par\noindent\textbf{Problem.}~}
\newcommand{\IRwhy}{\par\smallskip\noindent\textbf{Why it matters.}~}
\newcommand{\IRres}{\par\smallskip\noindent\textbf{Proposed resolution.}~}
\newcommand{\IRmeta}[4]{\par\smallskip\noindent\textbf{Owner:}~#1\enspace\textbar\enspace
  \textbf{Priority:}~\IRprio{#2}\enspace\textbar\enspace\textbf{Phase:}~#3\enspace\textbar\enspace
  \textbf{Status:}~#4}
\newcommand{\IRsrc}[1]{\par\noindent\textbf{Sources:}~#1}
\newcommand{\libc}[1]{\textsf{lib-L#1}}

%===============================================================================================
\subsection{How to read the register}\label{subsec:iss-read}

\paragraph{Priorities.} Each open issue has one of three priorities.
\begin{description}[nosep,leftmargin=2.6em,labelwidth=2em,font=\normalfont\bfseries]
\item[P0] \emph{Launch-blocking.} The issue must be closed before the contract is frozen and the
  first season (T0) opens. P0 issues change the trusted statement, or they are needed to show that
  the certificate is inhabited and checkable within budget.
\item[P1] \emph{Before T1.} The issue must be closed before the RISC-V track T1 opens. T0 can run
  without it.
\item[P2] \emph{Later.} The issue concerns T2, T3, a further track or a research question. It does
  not block any season before its own.
\end{description}

\paragraph{Owner types.} The \emph{contract team} writes and maintains the trusted Lean contract,
its controls and the machine. The \emph{library team} writes the untrusted, kernel-checked
organiser library and the reference submissions. The \emph{harness team} owns the checking
pipeline, the emulator, the sampler, the badge path and the benchmark generators. The
\emph{cryptography reviewers} check definitions, cited theorems, constants and hypotheses against
the literature. \emph{Process} means a decision or a rule that the organisers make and publish.

\paragraph{Phases.} Phases 0--5 are those of the roadmap (\cref{subsec:iss-roadmap}): Phase~0
contract freeze, Phase~1 T0 season, Phase~2 T1, Phase~3 T2, Phase~4 further tracks, Phase~5 T3.
\emph{Status} is one of \emph{open}, \emph{decision needed}, \emph{in progress} or \emph{partial}
(some of the work is done and measured).

\paragraph{Identifiers and sources.} Open issues are numbered by theme: C (definitions and
contract), M (machine and cost model), L (library and proof tooling), H (harness and badge path),
B (benchmark), R (recursion and T2), G (governance and process) and Q (research questions). Source
identifiers are those of the underlying documents:
\begin{itemize}[nosep]
\item A1-C1, A2-M3, A3-m4 are findings of reviews A1--A3, and A3-S3 is A3's simplification S3;
  A4-F1--F7 are A4's findings, A5-M$n$ A5's, A6-G0, P1, G1, G2, R1, H1, F1, T1 and N1--N6 A6's,
  and A7-M$n$ A7's;
\item I-BAS-1 and I-BAS-2 are the baseline's contract findings (\cref{subsec:base-findings});
\item D1-G5 and D2-G14 are the residual gaps of D1 §7 and D2 §8, and G-D3-$n$ those of D3 §7;
\item O8--O14 are the open items of DESIGN §14;
\item R1--R8 are the research notes, cited with their implications section;
\item U1--U3 are the three user inputs (PLAN §5);
\item W3 marks objections raised while this report's recursion, cost and benchmark sections were
  written, as recorded in DESIGN §14.
\end{itemize}

\paragraph{Two uses of \enquote{L}.} Security \emph{levels} L0--L4 are those of \cref{sec:protocol}
(L1 is the v1 contract, L2 adds the tree theorem and efficient extraction, and so on). The
components of the organiser library (A3 §4) are written \libc{0}--\libc{6} here to keep the two
apart. The critical path is \libc{4}, the knowledge-sound hash-based PCS over binary codes.

%===============================================================================================
\subsection{Open issues}\label{subsec:iss-open}

\Cref{tab:iss-overview} lists the open issues. The boxes below give each P0 and P1 issue in full,
except the five v0.4 contract items, which share \cref{tab:iss-v04}; \cref{tab:iss-p2} gives the P2
items. There are 47 open issues: 15 at P0, 23 at P1 and 9 at P2. Statuses reflect the compiled v0.3
contract with the round-3 and round-4 harness fixes (\code{NOTICE-P2d.md}, \code{NOTICE-P2e.md}) and
the baseline certified against it. The v0.4 contract items (C-5, M-5--M-8) are deliberately \emph{not}
in the freeze: the counters and the machine the baseline is certified against
stay fixed, and v0.4 is released with a new contract id at a season boundary before T1.

\begingroup\small
\renewcommand{\arraystretch}{1.08}
\begin{longtable}{@{}l>{\raggedright\arraybackslash}p{0.47\textwidth}>{\raggedright\arraybackslash}p{0.22\textwidth}cc@{}}
\caption{Open issues at a glance.}\label{tab:iss-overview}\\
\toprule
ID & Issue & Owner & Prio. & Phase \\
\midrule
\endfirsthead
\toprule
ID & Issue & Owner & Prio. & Phase \\
\midrule
\endhead
\bottomrule
\endlastfoot
C-1 & Release candidate: the v0.3 contract with its controls & contract & P0 & 0 \\
C-3 & Metaprogramming and compiled-code policy for submissions & process, contract & P0 & 0 \\
C-4 & \lean{QProg} certificate and T0 scoring & contract, harness & P0 & 0 \\
C-5 & Per-hart fuel rule in deployment (\lean{NoTimeout} or a global budget) & contract, process & P1 & 1 \\
C-6 & Interactive soundness union bound & library, crypto & P1 & 1 \\
C-7 & Computable is not feasible: indexer and prover feasibility rests on the badge path & process, harness & P1 & 1 \\
M-1 & Zbc in the machine, decoder and emulators & contract, harness & P0 & 0 \\
M-2 & Class set and headline aggregation & process, crypto & P0 & 0 \\
M-3 & Calibrate the judgement constants & crypto, harness, process & P1 & 1 \\
M-4 & Field-syscall (T1-F) and precompile board & contract, crypto & P2 & 4 \\
M-5 & \texttt{SPAWN}/\texttt{JOIN} pricing on the recursion board & contract & P1 & 1 \\
M-6 & Leaf hashing on the latency board & contract & P1 & 1 \\
M-7 & Timer structural hazards (memory port, divider, WAW) & contract, harness & P1 & 1 \\
M-8 & Power-of-two oracle lengths and leaf padding & contract & P1 & 1 \\
L-1 & \libc{0}: machine facts & library & P0 & 0 \\
L-2 & \libc{1}: symbolic executor and accepting-run cost judgement & library & P0 & 0 \\
L-4 & \libc{2}: verified field routines for $K$ and $E$ & library & P1 & 1 \\
L-5 & \libc{3}: IOP components on the local definitions & library & P1 & 1 \\
L-6 & \libc{4}: knowledge-sound PCS over binary RS codes (critical path) & library, crypto & P1 & 1--2 \\
L-7 & \libc{5}: reference succinct submission & library & P1 & 2 \\
L-8 & Bridge to ArkLib in both directions & library & P1 & 1 \\
L-9 & Executable $K$/$E$ runtime for the indexer and the honest prover & library, harness & P1 & 1 \\
L-10 & Hart combinators: race freedom, work/span composition, footprint & library & P1 & 1 \\
H-1 & CI pipeline and prebuilt trusted environment & harness & P0 & 0 \\
H-3 & Fiat--Shamir prefix, encodings and sampler: specification and test vectors & harness, crypto & P1 & 1 \\
H-4 & Trusted emulator & harness & P1 & 1 \\
H-5 & Badge path implementation & harness, process & P1 & 1 \\
H-6 & Image literal rendering and admission & harness & P1 & 1 \\
H-7 & Production sandbox: real landrun instead of the bubblewrap shim & harness & P0 & 0 \\
H-9 & Specified but unbuilt harness components (preflight, source policy, \code{elf2image}, \lean{SkipKernel} regression) & harness & P0 & 0 \\
B-1 & Benchmark generators over $K$ & harness & P1 & 1 \\
B-2 & Public-input packing and Booleanity convention & harness, crypto & P1 & 1 \\
R-2 & Tree composition theorem & library, crypto & P2 & 3 \\
R-3 & All-runs cost bound for deciders & contract & P2 & 3 \\
R-4 & Bounded-depth accumulation & contract, crypto & P2 & 4 \\
G-1 & Review of the release candidate and statement-bug bounty & crypto, process & P0 & 0 \\
G-2 & Contract freeze and pinned release & process, contract & P0 & 0 \\
G-3 & Versioning and record-revalidation machinery & harness, process & P0 & 0 \\
G-4 & Hypothesis governance and the first named hypothesis & crypto, process & P1 & 1 \\
G-5 & Kernel, comparator and toolchain upgrades & process, harness & P1 & 1 \\
G-6 & Second-kernel policy (nanoda, \code{leanchecker}) & process, harness & P0 & 0 \\
G-7 & Trusted-surface pins: unpinned binaries, host tools, elaborator and Lake & harness, process & P0 & 0 \\
Q-1 & Efficient extraction (security level L2) & crypto, contract & P2 & 4 \\
Q-2 & Recursive soundness after hash instantiation & crypto & P2 & 3--4 \\
Q-3 & Arithmetised recursion (T3) & crypto, contract & P2 & 5 \\
Q-4 & ROM argument track (security level L3) & contract, crypto & P2 & 4 \\
Q-5 & Further tracks: Boolean $\mathrm{GF}(2)$ and certified attacks & contract, crypto & P2 & 4 \\
\end{longtable}
\endgroup

%-----------------------------------------------------------------------------------------------
\subsubsection{Definitions and contract}

\begin{issue}{C-1}{Release candidate: the v0.3 contract with its controls}
\IRprob The v0.3 contract is compiled (\code{contract-v03/}): 12 modules and a root, 796 code lines,
no \lean{sorry}, with positive controls (toy protocol, \lean{fullRecAcc}, machine smoke run with the
18 counters, first touch and the timer, harts, an oracle-reading decider) and negative controls (A2-C1,
C3, M3, accept-all, reject-all, error cap, five hart-region violations), all in CI. It includes the
A4 fixes F1--F7, the A5 changes, the sparse \lean{Index}, reflecting padding and the round-3 fixes
(A6-G0, P1, G1, G2, R1). What is missing for a release candidate: D1's controls C1--C9 mapped one to
one to \cref{tab:sec-controls}; the T0 certificate (C-4); the A6 notes that touch the trusted text
(N2: \lean{Profile.w} as \lean{Fin 15 → ℚ} instead of a zipped list; N1: reworded gate
documentation); a published trusted-line count; and a CI step that checks the appendix listings
against the source (\path{report/scripts/regen-appendix-a.py --check}).
\IRwhy Statement-level holes are invisible to comparator by construction
(\cref{sec:threats}). All four attacks built in Lean against v0.1 were definitional. A contract that
has not been compiled together with its controls cannot be frozen.
\IRres
\begin{enumerate}[nosep,leftmargin=1.4em]
\item Freeze the contract at v0.3 plus the round-3 fixes. Keep \code{Machine.lean} and the 18
  counters as certified by the baseline; collect the v0.4 changes (C-5, M-5--M-8) separately.
\item \code{Controls/Negative.lean}: every negative control of \cref{tab:sec-controls} must fail to
  elaborate or be refuted by a proved theorem. A CI job fails if one of them elaborates.
\item \code{Controls/Positive.lean}: add C4, the certified baseline (V51).
\item Keep ArkLib and CompPoly outside the contract's import closure (\code{Bridge/} and the library
  only).
\item Publish the trusted-line count against the v0.3 target of at most 850 code lines excluding the
  vendored \code{Rv64/*} files (v0.3: 796; v0.2: 582).
\end{enumerate}
\IRmeta{contract team}{P0}{0}{partial (v0.3 compiled with controls; round-3 fixes in; release
  candidate open)}
\IRsrc{O8; A2 fix list 1--3, 8, 9; A2-m3; D1 §5, D1-G3; A3 fix list 1--2; A4; A6; contract
  \code{RESULT.md} §0, §R3; R1 (non-vacuity witness); R5 §7.1 (refutation lemmas as guard theorems).}
\end{issue}

\begin{issue}{C-3}{Metaprogramming and compiled-code policy for submissions}
\IRprob sig.golf's reflective executor relies on \lean{implemented\_by} and \lean{evalExpr}
metaprograms. The computability gate bans \lean{noncomputable}, \lean{implemented\_by},
\lean{extern} and \lean{csimp} in the closure of \lean{VK}, \lean{IDX}, \lean{Strategy} and the
extractor. The rule for the rest of a submission (proof-producing tactics, local macros, reflection)
is undecided (\cref{sec:threats}, attack L3).
\IRwhy The rule decides whether submitters may use the fastest proof route, and whether compiled
code can influence anything that the badge path evaluates. It must be published before the first
submission.
\IRres Publish a three-part rule. (a)~Metaprograms are allowed in the organiser library, because the
kernel re-checks every term they produce. (b)~In submissions they are allowed in proof-producing
code whose output the kernel re-checks, and never in the closure of a gated definition.
(c)~Environment commands (\lean{addDecl}, \lean{debug.*} options) stay refused by the
comment-stripped header audit. Add a negative control: a \lean{VK} with \lean{implemented\_by}
fails the gate.
In v0.3 the gate enforces the compiled-code half of the rule over \emph{every} untrusted module:
no \lean{implemented\_by}, \lean{extern}, \lean{csimp}, \lean{init}, \lean{builtin\_init} or
\lean{export} entry, so a submission may not use \lean{initialize}, \lean{register\_simp\_attr} or
\lean{@[export]} (A6-G1); debug options are caught only by the kernel replays (A6-N1). Since
round~4 the computability rule covers only the badge-path fields, so declared errors may use
\lean{ℝ≥0} operations (A6-N3, A7-M2), and a solution must elaborate deterministically, because the
gate tree's build must equal the comparator-verified one (A7-M3).
\IRmeta{process, contract team}{P0}{0}{partial (enforced by the gate; publication and the rule for
  proof-producing metaprograms pending)}
\IRsrc{A3-m4; A2-M4; A4-F3; A6-G1, N1, N3; R2 §6.1 (policy checker); R3 §6.4 (textual audit of
  metaprogramming).}
\end{issue}

\begin{issue}{C-4}{\lean{QProg} certificate and T0 scoring}
\IRprob T0 is the first season. Its verifier is a contract query program \lean{QProg} (a free monad
over \code{QUERY} with pure local computation), scored by virtual query cost and bytes. The type and
its evaluator are drafted (42 lines). Three parts are missing: the T0 certificate (completeness,
RBR-KS and the cost bound stated over \lean{QProg}), the T0 scorer, and the T0 admission caps.
\IRwhy Without them the Phase~1 season cannot open. The \lean{QProg} midpoint is also how T0
records carry over into T1.
\IRres Write \code{Certificate.T0}: decision agreement of the \lean{QProg} with the spec verifier,
and a virtual cost bound per class over all raw oracle contents. Use the same query price as the
machine, $\mathit{depth}_j(P_{H,P}+35)$ per call plus the leaf bound
$P_{H,P}(q^{\mathrm{bytes}}/64+q^{\mathrm{count}})$ (\cref{subsec:cost-zk}), plus bytes, so that T0
and T1 scores share their query component. Add a library lemma for \libc{1}: a machine image that
refines a \lean{QProg} inherits \lean{Agrees} and the query part of the cost. Give the T0 board
positive and negative controls (the baseline as a \lean{QProg}, and an accept-all \lean{QProg}).
\IRmeta{contract team, harness team (scorer)}{P0}{0}{in progress (\lean{QProg} type drafted)}
\IRsrc{O10; A3-M5; A3-S6; D2 §6.1; DESIGN §11.}
\end{issue}

\begin{issue}{C-6}{Interactive soundness union bound}
\IRprob The library proves the deterministic core of straight-line extraction
(\lean{no\_bad\_imp\_false}, \lean{chain\_extracts}), but not the probabilistic step: the accepting
fraction on a false statement is at most $\sum_c\eps_c$, by prefix independence of the challenge at
round $c$. The hit control is proved directly by counting. The gap is unchanged since v0.2.
\IRwhy Without it, the step from the contract's per-round \lean{RbrKS} to an interactive soundness
error rests on the cited theorems alone, and auditors cannot check the error accounting in Lean.
\IRres Prove the union bound in the library over the contract's round shapes and challenge menu
(sampling is exact over $K$ and $E$, so the per-round probability is a counting statement), and add
it as a positive control.
\IRmeta{library team, cryptography reviewers}{P1}{1}{open}
\IRsrc{contract v0.2 \code{RESULT.md} §8 item 1; contract v0.3 \code{RESULT.md} §10 item 6;
  \cref{subsec:sec-rbrks}.}
\end{issue}

\begin{issue}{C-7}{Computable is not feasible: indexer and prover feasibility rests on the badge path}
\IRprob Gate rule~4 (\code{--computable}) guarantees that the badge-path fields \lean{VK},
\lean{IDX}, \lean{prover} and \lean{extractor} have compiled code. It does not make them feasible: a
satisfiability table $x\mapsto[\exists w,\ \mathsf{R1CS}]$ is Lean-computable, by enumeration over the
finite $K^{n_w}$ with the contract's \lean{Decidable} instance for \lean{R1CS}. Lean's
\lean{noncomputable} marks definitions without compiled code (choice, real inverses), not expensive
ones. (Round~4 fixed the related A7-M2: rule~4 now checks only the badge-path fields, so a declared
error such as \lean{errorCap} no longer fails the gate; V54.)
\IRwhy Feasibility, and hence the closure of A2-M4 for records, rests entirely on the badge path's
time and memory budgets on hidden indices, which are not yet implemented (H-5); the rules must not
claim more for the gate.
\IRres State in the rules that the gate ensures executability and the badge path feasibility, for
records only; set the badge budgets from measured compiled indexers (L-9); add a negative control: a
brute-force indexer passes the gate and fails the badge budget.
\IRmeta{process, harness team}{P1}{1}{open (gate scope fixed in round~4; badge budgets and rules text open)}
\IRsrc{A7-M1, A7-M2; A6-N3; A2-M4; \code{NOTICE-P2e.md}; \cref{subsec:sec-extractor}; \cref{tab:thr-catalogue} row S8.}
\end{issue}

%-----------------------------------------------------------------------------------------------
\subsubsection{Machine and cost model}

\begin{issue}{M-1}{Zbc in the machine, decoder and emulators}
\IRprob The machine adds \code{clmul}, \code{clmulh} and \code{clmulr} to RV64IM. riscv-zkvm's
hand model has no Zbc, so the semantics, the decoder arms and the price class are owned by the
contract. The same holds for the W-ops. Neither is covered by riscv-zkvm's Sail-equivalence proofs.
The raw-word decoder is trusted and tested, not proved.
\IRwhy If the decoder or the semantics diverges from the hardware, every certified theorem is about
the wrong machine. Binary-field verifiers spend most of their cycles in \code{clmul}.
\IRres
\begin{enumerate}[nosep,leftmargin=1.4em]
\item Done: the three instructions are defined in \code{Machine.lean} (via \lean{K.clmul}). Add the
  library lemma relating them to CompPoly's \lean{carryLessMul} (A3-n1).
\item Run differential tests of decoder and semantics against \code{qemu-riscv64}, which supports
  Zbc \cite{RISCV-Bitmanip}, and against \code{cyclemu}, on the D3 kernels and on random words,
  extended to branches, \code{JAL}/\code{JALR}/\code{AUIPC} and fetch (not covered by A4's tests).
  Proposed acceptance: zero mismatches.
\item Run \code{check-vendor.sh} in CI.
\item State in the trusted-surface table which instructions are covered by Sail equivalence and
  which are contract-owned.
\end{enumerate}
\IRmeta{contract team (semantics), harness team (tests)}{P0}{0}{partial (semantics in the v0.3
  machine; A4 ran about 9{,}000 decoder and semantics cases against QEMU with 0 mismatches; the
  baseline diff-tested counters, first touch and timer against \lean{runFull} on two tiny instances)}
\IRsrc{G-D3-1; D2-G4; A3-n1; R6 §6.1; R8 §5.4; \cref{subsec:thr-residual} item 4.}
\end{issue}

\begin{issue}{M-2}{Class set and headline aggregation}
\IRprob At the worst-case $\mathsf{nnz}=16N$, all three classes are succinct and nearly flat for
the modelled designs, so $\mathrm{Score}_{\mathrm{T1}}$ is about $C_{20}$. A design whose cost grows
faster in $k$ can win the geometric mean and lose at $2^{24}$, which is where recursion lives. A1
proposed the classes $\{20,22,24\}$, or $C_{24}$ as the headline. v0.2 keeps $\{16,20,24\}$ and
does not record why.
\IRwhy Class shapes are part of the certificate. Changing them after the freeze invalidates every
proof.
\IRres Decide before the freeze, and record the rationale in DESIGN. Recommendation: keep
$\{16,20,24\}$ in the certificate, because $2^{16}$ keeps the baseline and hybrid designs cheap to
demonstrate. Keep the geometric mean as the score. Publish per-class boards with $C_{24}$ displayed
next to the score. State in the rules that the aggregation is a scorer-side function of the
certified vectors and may be revised at a season boundary without re-proof.
\IRmeta{process, cryptography reviewers}{P0}{0}{decision needed}
\IRsrc{A1-m2; A1-n3 (per-class boards); D2-G2.}
\end{issue}

\begin{issue}{M-3}{Calibrate the judgement constants}
\IRprob Several constants are judgements (\cref{subsec:cost-fails,subsec:cost-sensitivity}):
\begin{itemize}[nosep]
\item the per-profile weights and hash prices $P_{H,P}$ (RISC~Zero and Jolt are guesses) and $+35$
  per level;
\item the per-profile FS coefficients $2P_H+30$, $0.75P_H+6$ and 4;
\item the blow-up threshold 16 and exponent $1/4$, which were tuned on prime-field SPARK, whereas the
  binary model gives $V\approx30$ at $\rho=1/2$ and rate $1/16$ only ties rate $1/4$;
\item the decider amortisation $/64$.
\end{itemize}
The model itself is uncertain to about a factor of 2 on succinct rows (SPARK stack size,
weight evaluations per query, fold schedule).
\IRwhy The constants decide whether hash-heavy or arithmetic-heavy designs win. Re-weighting needs
no re-proof, but changing the scorer during a season reorders the board.
\IRres
\begin{enumerate}[nosep,leftmargin=1.4em]
\item Set the FS constants by measuring a BLAKE2s-based Fiat--Shamir transcript on the probe
  emulator.
\item Collect SP1, OpenVM and ZisK row counts per compression, per multiplication and per memory
  access; publish the derivation of $P_H$, $+35$ and the MUL weight.
\item Instrument leanVM's WHIR verifier path to validate the trajectory model \cite{leanVM}.
\item Consider counting only committed (level-0) oracles in $V_k$.
\item Version the constants with the scorer. Freeze them per season, and change them only at
  season boundaries, publishing the board under both the old and the new constants.
\item Freeze $\lambda=1/8$, and do not advertise it as a size incentive.
\end{enumerate}
\IRmeta{cryptography reviewers, harness team (measurements), process (versioning)}{P1}{1}{open}
\IRsrc{O12 (W3 objection 4); D2-G1, D2-G2; G-D3-5, G-D3-6, G-D3-10; A1-M5, A1-m1;
  \cref{tab:thr-catalogue} row C14.}
\end{issue}

\paragraph{The v0.4 contract items.} Five P1 items change the trusted machine or its pricing and
are documented gaps of v0.3, deliberately not in the freeze (\cref{tab:iss-v04}). Each is a
score-meaning or deployment gap, not a soundness hole of the scored theorem. They land together in
v0.4 under a new contract id before T1 (owner: contract team; phase~1; status: open).

\begingroup\small
\renewcommand{\arraystretch}{1.08}
\begin{longtable}{@{}l>{\raggedright\arraybackslash}p{0.40\textwidth}>{\raggedright\arraybackslash}p{0.36\textwidth}>{\raggedright\arraybackslash}p{0.10\textwidth}@{}}
\caption{The v0.4 contract items (P1).}\label{tab:iss-v04}\\
\toprule
ID & Gap & Proposed resolution & Sources \\
\midrule
\endfirsthead
\toprule
ID & Gap & Proposed resolution & Sources \\
\midrule
\endhead
\bottomrule
\endlastfoot
C-5 & \emph{Per-hart fuel rule in deployment.} A child runs with its parent's remaining fuel and the parent is not charged, so a run may take $(\#\mathrm{spawns}+1)\cdot2^{36}$ steps; a \enquote{loop past the fuel, then \texttt{HALT(1)}} backdoor is rejected by the model but accepted by a deployment that does not copy this rule (a native verifier counts no instructions) & a certificate clause \lean{NoTimeout} (every input halts or faults before the fuel runs out; proof burden on rejecting paths), or one global step budget for all harts stated in the challenge text & A6-F1, row H12 \\
M-5 & \emph{\texttt{SPAWN}/\texttt{JOIN} on the recursion board} count one $n_{\mathrm{sys}}$ each at weight~1, but a zkVM image needs a context switch (about 31 stores and 31 loads): a register-preserving call for about 3 units & a counter $n_{\mathrm{spawn}}$ priced at $64(w_{\mathrm{ld}}+w_{\mathrm{st}})/2$ (ZisK $\approx38$, SP1 $\approx67$, RISC~Zero and Jolt 64), or $n_{\mathrm{sys}}$ weight 32 in every profile & A6-H1, row C26 \\
M-6 & \emph{Leaf hashing on the latency board.} The timer's \texttt{QUERY} bubble charges $\mathtt{QUERY\_LEVEL\_CYCLES}\cdot\mathsf{depth}+2^{\mathit{logLeaf}}$, one cycle per leaf element, about 195 times below its hash cost & charge $\mathtt{QUERY\_LEVEL\_CYCLES}\cdot(\mathsf{depth}+\lceil2^{\mathit{logLeaf}}/8\rceil)$ in the timer & A6-P1, \cref{subsec:cost-lat} \\
M-7 & \emph{Timer structural hazards:} two load/store slots per cycle, a fully pipelined 16-cycle divider and no WAW stall; memory-bound loops look about 2 times too fast & one memory port and an unpipelined divider (about 6 trusted lines), re-proving the composition lemmas; or calibrate on the reference board & A6-T1, row C27 \\
M-8 & \emph{Power-of-two oracle lengths.} The baseline's index needs $35\cdot2^k$ of its $2^{k+6}$ words; the padding is about 39\% of its RISC~Zero score (13\% under v0.2), and its message carries 1{,}025 ignored entries & allow \lean{OShape} lengths that are multiples of the leaf, with $\mathsf{depth}=\lceil\log_2(\ell/\mathsf{leaf})\rceil$ for pricing & I-BAS-2 \\
\end{longtable}
\endgroup

%-----------------------------------------------------------------------------------------------
\subsubsection{Library and proof tooling}

The organiser library (\cref{tab:iss-lib}) is untrusted, kernel-checked and importable. \libc{0}
and \libc{1} are the starter tooling (the baseline's \code{Rv/L0.lean} is a \libc{0} prototype).
\libc{2} and \libc{3} gate any succinct entry. \libc{4} followed by \libc{5} is the critical path to a
competitive leaderboard. A3 estimates the elapsed effort at 3--6 organiser agent-weeks for
\libc{0}--\libc{2}.

\begin{table}[h]
\centering\small
\begin{tabularx}{\linewidth}{@{}l>{\raggedright\arraybackslash}Xrll@{}}
\toprule
Comp. & Content & Lines [A3] & Needed for & Issue \\
\midrule
\libc{0} & machine facts: step, decode, access, readback, layout lemmas & 1.5--2.5k & every T1 entry & L-1 \\
\libc{1} & symbolic executor, loop combinators, \lean{QProg} segments, accepting-run cost judgement & 3--5k & every T1 entry & L-2 \\
\libc{2} & verified $K$/$E$ routines with specs and cycle counts & 2--4k & T1 succinct & L-4 \\
\libc{3} & sumcheck, batching, GKR, query check, append, frame rule on the local definitions & 3--7k & T0/T1 succinct & L-5 \\
\libc{4} & knowledge-sound PCS over binary RS codes (interim: tensor-code PCS) & 15--40k (5--10k) & any succinct entry & L-6 \\
\libc{5} & reference succinct submission & 10--20k & viability, seeding & L-7 \\
\libc{6} & harness: renderer, sampler, emulator, prover sandbox, generators & 4.5--6.5k & launch & H-1--H-6, B-1 \\
\bottomrule
\end{tabularx}
\caption{Organiser library plan (A3 §4).}\label{tab:iss-lib}
\end{table}

\begin{issue}{L-1}{\libc{0}: machine facts}
\IRprob There is no importable RISC-V proof tooling for the contract machine. Stock tactics over
the hand model hit measured traps; one costs 266\,s of kernel time and then fails on a single
\code{JAL} step.
\IRwhy \libc{0} gates every T1 entry. The baseline needed 1{,}920 lines without a library (A3 had
estimated 2.7--4.5k) because its image has 84 instructions; succinct entries, with thousands, are
where the library matters.
\IRres Prove symbolic-operand step lemmas for every contract instruction, including the W-ops and
Zbc. Add raw-word decode lemmas, \lean{accessValid} lemmas, a closed-form \code{QUERY} and tape
readback, and layout facts (1.5--2.5k lines). Avoid the documented traps: use lemmas with
symbolic immediates, no \lean{decide} on \lean{BitVec}, and no \lean{n\#w} literals.
\IRmeta{library team}{P0}{0}{open}
\IRsrc{A3-C3, A3 §4; D2-G5; D1-G11; R6 §6 item 2.}
\end{issue}

\begin{issue}{L-2}{\libc{1}: symbolic executor and accepting-run cost judgement}
\IRprob Proofs about loop-based verifiers need an executor and loop rules. Unrolling does not
scale: \lean{runBlock} is superlinear, at 161\,s for a 128-instruction block. ots.golf records
already take 14--25\,min to check.
\IRwhy Check time is the binding constraint. Proofs must scale with invariants, not with the
number of instructions executed.
\IRres Port sig.golf's \lean{symRun}, \lean{symRun\_sound}, \lean{BlockSpec} and
\lean{Steps.iterate} to the contract machine. Add counted-loop and early-exit-loop combinators (the
latter of the form \enquote{$\forall$ rows ok $\leftrightarrow$ no reject}), \lean{QProg}-segment
combinators (C-4), and the accepting-run cost judgement (sig.golf's \lean{Good}). Keep blocks small.
3--5k lines.
\IRmeta{library team}{P0}{0}{open}
\IRsrc{A3-C3, A3 §4, A3-M5; D2-G5; R1 (check time); R2 §6.2; R6 §6 item 2.}
\end{issue}

\begin{issue}{L-4}{\libc{2}: verified field routines for \texorpdfstring{$K$}{K} and \texorpdfstring{$E$}{E}}
\IRprob Refinement proofs from \lean{BitVec} to field arithmetic are the hardest part of a T1
submission's machine-level proof, and every entrant would otherwise repeat them.
\IRwhy DESIGN §11 opens T1 when \libc{0}--\libc{2} ship. Without shared routines, assembly
engineering dominates the score.
\IRres Ship verified RISC-V routines with specifications and cycle counts:
\begin{itemize}[nosep]
\item $K$ multiplication (\code{clmul} plus two folds);
\item lazy-XOR accumulation followed by one reduction;
\item $E=K[y]/(y^3+y+1)$ multiplication and squaring;
\item inversion checked against a prover-supplied hint;
\item $\mathrm{eq}(r,\cdot)$ and Lagrange evaluation loops;
\item \lean{rfl}-bridges to CompPoly's \lean{BF64} and \lean{BF64.Ext3}.
\end{itemize}
2--4k lines.
\IRmeta{library team}{P1}{1}{open (field-level facts exist in \code{Library/Field.lean})}
\IRsrc{G-D3-4; A1-M5c; A3 §4, A3-S4; R8 §5; DESIGN §2.}
\end{issue}

\begin{issue}{L-5}{\libc{3}: IOP components on the local definitions}
\IRprob ArkLib's composition theorems (\lean{append}, \lean{seqCompose}, Fiat--Shamir) and its
sumcheck RBR-KS are not sorry-free \cite{arklib2026}. leanerVM has sorry-free sumcheck, batching,
GKR grand-product and query-check components, but they are stated over ArkLib's definitions, not
over the contract's local \lean{RbrKS}/\lean{KSF}.
\IRwhy Every succinct T0 or T1 entry needs these components. Without a library, each submission
would have to prove sumcheck from scratch.
\IRres Port or bridge leanerVM's components to the local definitions. Add sequential composition
(append), the frame rule (D1-G2), and \lean{Refinement.knowledge\_transport} for moving
knowledge soundness along witness maps. 3--7k lines; upstream to ArkLib where possible.
\IRmeta{library team}{P1}{1}{open}
\IRsrc{A3 §4; R5 §7.1; R4 §9 items 2, 7; D1-G2.}
\end{issue}

\begin{issue}{L-6}{\libc{4}: knowledge-sound PCS over binary RS codes (critical path)}
\IRprob A competitive entry needs a hash-based polynomial commitment with RBR \emph{knowledge}
soundness over Reed--Solomon codes on additive domains in $K$, with challenges in $E$. No such
formalisation exists:
\begin{itemize}[nosep]
\item ArkLib has no WHIR, and its FRI requires \lean{NonBinaryField};
\item the FRI-Binius and ring-switching theorems are sorried;
\item leanerVM targets WHIR at RBR \emph{soundness} only, and has not built it;
\item no correlated-agreement theorem for a subfield code under extension-field challenges is
  formalised.
\end{itemize}
\IRwhy A3 estimates that a competitive entry costs 45--100k lines with the organiser library and
100--170k without it. Without \libc{4} and \libc{5}, the board can be expected to stay at the
trivial or hybrid baseline for months.
\IRres Work in two stages and staff \libc{4} first.
\begin{enumerate}[nosep,leftmargin=1.4em]
\item \emph{Interim} (5--10k lines): a tensor-code PCS in the style of Ligero and Brakedown
  \cite{AHIV17-Ligero,GLSTW23-Brakedown}, from the proved tensor and interleaved proximity gaps
  \cite{DG25-Interleaved}. It runs as a flagged track (A3-S8).
\item \emph{Full} (15--40k lines): a unique-decoding, knowledge-sound RS PCS over additive domains,
  from ArkLib's sorry-free UD correlated agreement \cite{BCIKS20}.
\end{enumerate}
Either prove the extension-challenge correlated-agreement lemma, or restrict the PCS to codes over
$E$. Hold a risk review at the midpoint, and upstream the work to ArkLib.
\IRmeta{library team, cryptography reviewers}{P1 (critical path)}{1--2}{open}
\IRsrc{A3-M2, A3 §4, A3-S8; R8 §5 item 5; R4 §9 item 7; R3 §6 item 3; \cref{sec:benchmark}.}
\end{issue}

\begin{issue}{L-7}{\libc{5}: reference succinct submission}
\IRprob It is not known whether a succinct design can be certified within the T1 budget of 3\,h,
8 CPUs and 48\,GiB.
\IRwhy The reference submission proves viability, calibrates the budgets on a real succinct
proof, and seeds the board with a non-trivial record.
\IRres Build a Spartan- or GKR-style verifier \cite{Set20-Spartan} on \libc{2}--\libc{4} and
certify it on T1 (10--20k lines), first on the interim PCS and then on the full one. Publish its
check-time profile.
\emph{Exit for Phase~2:} certified within budget; badge path passed (prover 1\,h at $2^{20}$ and
8\,h at $2^{24}$); $\mathrm{Score}_{zk}$ below the certified baseline's $1.624\times10^9$
(\cref{tab:base-score}).
\IRmeta{library team}{P1}{2}{open (depends on L-4, L-5, L-6)}
\IRsrc{A3-M2, A3 §4, A3-M4; R2 §6.3.}
\end{issue}

\begin{issue}{L-8}{Bridge to ArkLib in both directions}
\IRprob The bridge from the local definitions to ArkLib's worst-case RBR-KS is proved. The
converse direction (ArkLib $\Rightarrow$ local) and the equivalence of the two completeness
notions (a \lean{perfectCompleteness\_of\_run\_support} argument) are not.
\IRwhy Agents who prove in ArkLib need the converse. It is also the sanity check that the local
restatement is not weaker than ArkLib's.
\IRres Write \code{Bridge/ToArkLib.lean}, reusing the lemmas of the existing direction, and complete
positive control C8. D1 estimates about one day.
\IRmeta{library team}{P1}{1}{partial (local $\Rightarrow$ ArkLib proved; still one-way in v0.3)}
\IRsrc{D1-G1; D1 §2.2; A2-m2; R4 §9 item 4.}
\end{issue}

\begin{issue}{L-9}{Executable \texorpdfstring{$K$/$E$}{K/E} runtime for the indexer and the honest prover}
\IRprob The badge path evaluates the compiled Lean indexer on hidden indices. Up to $2^{16}$ it also
compares the native prover's transcript with that of the compiled Lean prover. Lean-side field
arithmetic is slow:
\begin{itemize}[nosep]
\item compiled, about $28\,\mu$s per $K$ multiplication;
\item interpreted, about 1\,ms per $E$ multiplication;
\item $E$ inversion is not kernel-evaluable.
\end{itemize}
No fast executable implementation exists, and \lean{SampleableType} instances for $K$ and $E$ are
missing.
\IRwhy If the compiled indexer cannot preprocess a $2^{24}$ index within the badge budget, the
VK/IDX recomputation that closes A1-M3 cannot run, and no record can be granted at that class.
\IRres Ship an executable $K$/$E$ in the organiser library, on \lean{UInt64} with table- or
\lean{BitVec}-based carry-less multiplication, proved equal to the contract's $K$ and $E$. Benchmark
the compiled indexers of the baseline and of \libc{5} at all three classes. Set the badge budgets
from those measurements.
\IRmeta{library team, harness team}{P1}{1}{open}
\IRsrc{R8 §5 items 3 and 5; R5 §7.5; A2-M4; A1-M3; DESIGN §8.}
\end{issue}

\begin{issue}{L-10}{Hart combinators: race freedom, work/span composition, footprint}
\IRprob The library has no combinator for proving that a hart program respects its declared regions
(D4 estimates about 200 lines), no work/span composition lemma across \texttt{SPAWN}/\texttt{JOIN}
(the timer library composes straight-line blocks only), and only per-touch lemmas for first-touch
counts, no footprint combinator. The claim that every real interleaving reaches the sequential
semantics' state at \texttt{JOIN}, which is the latency board's deployment basis, is argued only in
prose (A6-N5); A6 found no counterexample.
\IRwhy Without these, every latency-board entry re-proves hart bookkeeping, and the deployment
argument for parallel runtimes cannot be audited in Lean.
\IRres Ship a region-discipline combinator, a W/S fold lemma over series-parallel spawn trees, a
footprint combinator for loops writing $[\mathit{base},\mathit{base}+8n)$, and either a written
proof or an interleaving-semantics lemma for the \texttt{JOIN} confluence claim.
\IRmeta{library team}{P1}{1}{open}
\IRsrc{contract \code{RESULT.md} §10 items 4--5; A6-N5; D4 §3, §5.4.}
\end{issue}

%-----------------------------------------------------------------------------------------------
\subsubsection{Harness and badge path}

\begin{issue}{H-1}{CI pipeline and prebuilt trusted environment}
\IRprob The organiser side of the pipeline of \cref{subsec:arch-ci} is assembled for v0.3:
\code{harness.sh} (stage order of Decision~A-3), \code{org-build.sh} (read-only organiser package),
\code{olean-manifest.py} (35,093-entry SHA-256 manifest), the gate, the renderer and \code{ci.sh}
(614\,s). Still missing: the submission-side intake (preflight, export at blob level and source
policy, H-9; verdict comments, record snapshots), the container image of the trusted build closure, and an isolated
packages setup; ArkLib, VCVio and CompPoly have no Mathlib-style cache, so a cold setup builds them
from source.
\IRwhy The pipeline is required for the T0 launch, and competitors must be able to reproduce
verdicts.
\IRres
\begin{enumerate}[nosep,leftmargin=1.4em]
\item Port the ots.golf and sig.golf pipeline (about 1k lines of Python and shell): sandbox,
  exact-module import allow-list, renderer, comparator, verdict comments and record snapshots.
\item Remove \code{packagesDir} from the lakefile and manifest.
\item Publish a container image of the trusted build closure, with its digest pinned under the
  contract id.
\end{enumerate}
\emph{Exit:} a clean-room rebuild from the published image reproduces the contract id and the
baseline's verdict.
\IRmeta{harness team}{P0}{0}{partial (organiser harness, manifest and CI built; intake and
  container open)}
\IRsrc{A3 §4 (\libc{6}); R1 implications; R2 §6.1; R3 §6 items 4--5; P1 probe §1.1; A4-F7;
  A6-G0; contract \code{RESULT.md} §R3.}
\end{issue}

\begin{issue}{H-3}{Fiat--Shamir prefix, encodings and sampler: specification and test vectors}
\IRprob Three pieces of the Fiat--Shamir step are harness code:
\begin{itemize}[nosep]
\item the transcript prefix
  \[H(\mathit{domSep}\,\|\,\mathit{contractId}\,\|\,\mathit{shapeHash}\,\|\,\mathit{claimHash}\,\|\,k\,\|\,\mathrm{MTroot}(\idx\,i)\,\|\,\vk\,i\,\|\,x);\]
\item the canonical encodings;
\item the masking of index challenges.
\end{itemize}
No written document yet fixes the hash, the Merkle layout of $\idx\,i$ or the domain separator.
\IRwhy A bug here is a real attack that comes with a valid certificate. Statement binding is the
mitigation against practical Fiat--Shamir attacks on generated R1CS \cite{KRS25,Fen26-GenR1CS}.
\IRres
\begin{enumerate}[nosep,leftmargin=1.4em]
\item Write an AUDIT document that specifies the transcript byte by byte.
\item Ship a reference implementation (sampler about 0.15k lines, plus the Merkle and Fiat--Shamir
  code) with published test vectors.
\item Optionally, state the sampler and the encodings as Lean definitions. Uniformity over $K$ and
  $E$ is then a one-line lemma, since every 8- or 24-byte string is an element.
\end{enumerate}
\IRmeta{harness team, cryptography reviewers}{P1}{1}{open}
\IRsrc{D1-G10, D1 §3.1; A2-m5 and A2's auditor checklist item 8; R7 §8 item 2; R8 §5 item 2.}
\end{issue}

\begin{issue}{H-4}{Trusted emulator}
\IRprob Demonstration runs need a trusted Lean emulator of the contract machine with sparse
$2^{36}$-byte memory and Zbc. A naive emulator runs at about $9\times10^6$ instructions per
second. At that rate the $2^{24}$ baseline run ($3.09\times10^9$ steps) takes about 6\,min; the
baseline's C emulator \code{rfemu} does it in 347\,s and 4.7\,GB (\cref{tab:base-emu}).
\IRwhy The check $\text{measured}\le\text{claimed}$ is the last line of defence against a
divergence between proof and machine. A failure halts the competition.
\IRres Write an optimised \lean{UInt64} emulator, targeting $3$--$10\times10^7$ instructions per
second (about 0.6k Lean lines plus 0.3k lines of tests). Either refine it to \lean{execute}, or
differentially test it against \lean{execute} on sampled traces and against \code{cyclemu} and
QEMU. Offer agents an untrusted fast path.
\emph{Exit:} zero mismatches; the $2^{24}$ baseline run takes at most 15\,min.
\IRmeta{harness team}{P1}{1}{open (an unverified C emulator exists in the baseline)}
\IRsrc{D2-G11; A3-m3, A3-M1, A3 §6.3; R6 §6 item 3; G-D3-12; A7-M4.}
\end{issue}

\begin{issue}{H-5}{Badge path implementation}
\IRprob The badge path of DESIGN §8 has several steps:
\begin{itemize}[nosep]
\item evaluate the compiled Lean indexer on the hidden indices and require a match with the native
  indexer;
\item run the sandboxed native prover against harness-sampled challenges over round IPC;
\item compare transcripts with the compiled Lean prover up to $2^{16}$;
\item check $\text{measured}\le\text{claimed}$;
\item enforce prover memory and wall-clock budgets (1\,h at $2^{20}$, 8\,h at $2^{24}$).
\end{itemize}
None of these is implemented; the v0.3 harness has a \code{--badge} hook that runs a command as its own
sandboxed process after the gate. The budgets are machine-specific, and gating by the native prover is
a practical rule, not a theorem.
\IRwhy The badge path is what closes the VK/IDX backdoor channel (A1-M3) and the
non-computable-indexer channel (A2-M4) for records.
\IRres Implement the steps above (native-prover sandbox, round IPC, Merkle and Fiat--Shamir code:
about 1k lines). Publish the reference machine. Define an appeal procedure for budget failures.
State in the rules that a badge failure affects records, not scores.
\IRmeta{harness team, process}{P1}{1}{open}
\IRsrc{DESIGN §8; A1-M3; A2-M4; A3-m2; D2-G9, D2 §2.6; D1-G6; \cref{subsec:arch-badge}.}
\end{issue}

\begin{issue}{H-6}{Image size check and lookup cost}
\IRprob The renderer chunks the image (256 words per literal), so a rendered 1\,MiB challenge
elaborates in 13\,s and 2.0\,GB. What remains is kernel evaluation: \lean{byteSize} $\le$ cap costs
about 70--80\,s once, and a naive lookup walks the prefix at 80\,s and 9\,GB per word, which
\lean{getElem?\_chunks'} avoids after a 28\,s length check.
\IRwhy These peaks can exceed the per-process memory limit and fail an otherwise valid T1
submission.
\IRres Ship the chunk lemmas as the documented route, render the image as its own module
(\cref{subsec:arch-budgets}), count the one-time checks against the budget, and re-measure under the
final pin.
\IRmeta{harness team}{P1}{1}{partial (chunked rendering and \lean{getElem?\_chunks'} done)}
\IRsrc{A3-m1; D2-G12.}
\end{issue}

\begin{issue}{H-7}{Production sandbox: real landrun instead of the bubblewrap shim}
\IRprob The v0.3 harness runs comparator and every solution build under
\code{scripts/bwrap-landrun.sh}, a bubblewrap shim with landrun's command line (read-only \code{/},
only the fresh tree's \code{.lake} writable, no network, cleared environment, own PID namespace).
The host has no landrun; landrun's executability restrictions (\code{--rox}) are not modelled, and
the packages directory is a shared read-only copy of the frozen v0.2 tree, not an isolated clean
setup.
\IRwhy The round-3 defences (Decision~A-3) assume the trusted artefacts are read-only to every stage
that runs solution code; the sandbox is what makes them so in production.
\IRres Run the harness under real landrun (or keep bubblewrap and publish it as the production
sandbox, with the same guarantees written down), with an isolated packages directory from a clean
cache setup, and re-run the \code{cmpproj} and \code{gateproj} regressions under it.
\emph{Exit:} both regressions are rejected under the production sandbox and order; the manifest
check passes on a clean-room organiser build.
\IRmeta{harness team}{P0}{0}{partial (bubblewrap shim in CI)}
\IRsrc{A4-F7; A6-G0; contract \code{RESULT.md} §R3.1, §10 items 1--2; \cref{subsec:arch-ci}.}
\end{issue}

\begin{issue}{H-9}{Specified but unbuilt harness components}
\IRprob \code{scripts/harness.sh} implements stages 0--2 (manifest, comparator first, gate) and a
\code{--badge} hook that runs a command as its own sandboxed process. The following are specified in
this report but not implemented: the submission preflight (systemd unit, seccomp, Landlock, refusing
root), the source and import allow-list and blob-level export (steps 1--3 of \cref{subsec:arch-ci});
the \code{elf2image} script offered to submitters; the badge checks (H-5); the organiser's trusted
emulator (H-4; only the baseline's unverified \code{rfemu} exists); and the Fiat--Shamir prefix and
sampler (H-3). \code{leanchecker} is not a harness stage (it runs in CI and in the baseline's
check), and the \lean{SkipKernel} probe has never run through comparator or the harness, so its
rejection rests on \code{leanchecker} evidence alone.
\IRwhy Launch needs the intake steps, and the trust story of \cref{sec:architecture} assumes them.
\IRres Implement steps 1--3 with the source policy of C-3, and \code{elf2image}; add \lean{SkipKernel}
to the harness regression; add a probe that compiled code cannot diverge from the kernel definition
without an attribute (A7-m17), or let the badge recompile the indexer in a clean process; keep an
implemented/planned status column in the CI step list and the trusted-surface table.
\IRmeta{harness team}{P0}{0}{open}
\IRsrc{A7-M4, A7-m17; A4-F7; contract \code{RESULT.md} §R3.1; \cref{subsec:arch-ci,subsec:arch-badge}.}
\end{issue}

%-----------------------------------------------------------------------------------------------
\subsubsection{Benchmark}

\begin{issue}{B-1}{Benchmark generators over \texorpdfstring{$K$}{K}}
\IRprob The benchmark has three families:
\begin{itemize}[nosep]
\item BLAKE2s compression as Boolean R1CS over $K$ with explicit XOR wires (40,992 rows), in
  chains of length 1, 24 and 399;
\item SHA-256 compression in the same encoding (59,704 rows), in chains of length 1, 16 and 268;
\item random sparse satisfiable R1CS with hidden seeds.
\end{itemize}
The circuits are counted symbolically, but there are no production generators with witness
generation.
\IRwhy A wrong or unsatisfiable generator silently breaks the demonstration runs and the badge
path.
\IRres Implement the generators with a reference witness generator, reusing leanVM's binary-field
R1CS where possible. Cross-check satisfiability against \code{baseline\_k.S} and against the
certified baseline. Publish the generator hashes. Check the rule that no benchmark index is
provably always satisfiable: each family must also produce unsatisfiable instances.
\IRmeta{harness team}{P1}{1}{open}
\IRsrc{D2-G10; G-D3-8; D1-G12; A3 §4 (\libc{6}); \cref{subsec:bench-families}.}
\end{issue}

\begin{issue}{B-2}{Public-input packing and Booleanity convention}
\IRprob The hash-chain public input must be packed so that $n_x\le2^{10}$, and the free input bits
of the first block need Booleanity rows. The convention is not fixed.
\IRwhy Without a fixed convention, the hash-chain instances either break the $n_x$ cap or leave
input bits unconstrained.
\IRres Define $x$ as 40 packed words, and add $w\cdot w=w$ rows for the first block's free input
bits. Fix this in the generator, and document it in the benchmark specification.
\IRmeta{harness team, cryptography reviewers}{P1}{1}{open}
\IRsrc{G-D3-13.}
\end{issue}

%-----------------------------------------------------------------------------------------------
\subsubsection{Governance and process}

\begin{issue}{G-1}{Review of the release candidate and statement-bug bounty}
\IRprob Three review rounds have run. Round~1 (A1--A3) found three critical holes in the submission
interface, each closed by a cheap definition change. Round~2 (A4) reviewed the compiled v0.2 contract
and found seven (F1--F7), all fixed in v0.3. Round~3 (A6) reviewed the compiled v0.3 contract and its
harness: no new statement-level hole, but a critical pipeline hole (G0), a high score-meaning hole
(P1), two gate bypasses and a renderer leak, all fixed (\cref{subsec:thr-evidence}). Each round found a new
class of hole at a lower layer. The bounty has not run.
\IRwhy Statement-level holes are the dangerous class, and no checker can find them; pipeline holes,
as G0 showed, can make a false statement pass every Lean check.
\IRres Freeze only a release candidate that has had one more review round with the pipeline in scope
(the production sandbox of H-7 and the harness order of Decision~A-3), from a cryptographer's and a
formal-methods perspective. Give every finding a written disposition. Run a public bounty for
statement and pipeline bugs between the release candidate and the freeze.
\emph{Exit:} no open critical or high finding.
\IRmeta{cryptography reviewers, process}{P0}{0}{partial (A4 and A6 done and dispositioned; final round
  and bounty open)}
\IRsrc{O11; A4; A6; \cref{subsec:thr-residual} item 10; \cref{subsec:thr-evidence}; PLAN §2 phase D;
  R5 §7.1 (disposition tables).}
\end{issue}

\begin{issue}{G-2}{Contract freeze and pinned release}
\IRprob The pins of DESIGN §13 are:
\begin{itemize}[nosep]
\item Lean and Mathlib v4.34.1;
\item ArkLib \repo{7653a90}, CompPoly \repo{572f997} and VCVio \repo{a4232d0};
\item riscv-zkvm \repo{6810fc1}, vendored byte-identically.
\end{itemize}
They are not yet frozen in a release. The P1 probe ran on v4.34.0 with ArkLib \repo{db7005f}; the
contract tree now uses v4.34.1. The prior competitions show the cost of churn: ots.golf made 46 pin
changes in 10 days, and sig.golf's accepted submission carries a 1.8k-line legacy-contract transfer.
\IRwhy Every contract change after launch costs re-proofs, record ports and trust.
\IRres
\begin{enumerate}[nosep,leftmargin=1.4em]
\item Check that \code{lake-manifest.json} matches DESIGN §13.
\item Freeze when the Phase~0 go/no-go criteria hold (\cref{subsec:iss-roadmap}), and publish the
  contract id with its changelog.
\item After the freeze, allow only security fixes. Each fix gets a new id and triggers revalidation
  (G-3).
\end{enumerate}
\IRmeta{process, contract team}{P0}{0}{open}
\IRsrc{D2-G14; R1 implications; R2 §6.3; R3 §6 item 6; A3-M6, A3 fix list 1, 3; DESIGN §13.}
\end{issue}

\begin{issue}{G-3}{Versioning and record-revalidation machinery}
\IRprob \Cref{subsec:arch-versioning} specifies contract ids, ported records with their provenance,
and three revalidation triggers (toolchain bump, machine change, hypothesis change). None of this
is implemented.
\IRwhy The machinery must exist before the first record. Retrofitting it is what made ots.golf's
post-launch churn expensive.
\IRres Store in every verdict the contract id and the checker commit. Keep a
\code{record-revalidations.json} ledger. Write a script that re-checks every affected record on a
trigger, and a porting procedure that keeps provenance and does not transfer credit.
\IRmeta{harness team, process}{P0}{0}{open}
\IRsrc{D2-G14; R1 implications; R7 §8 item 6.}
\end{issue}

\begin{issue}{G-4}{Hypothesis governance and the first named hypothesis}
\IRprob The conditional track (\cref{subsec:sec-hypotheses}) needs a closed, versioned
\code{Hypotheses.lean} and a written policy. Its first entry, leanerVM's \lean{McaJohnson}
transcription of \cite[Thm.~4.6]{BCHKS25} with the regime guard $\gamma\le1-\sqrt\rho-\eta$,
$\eta>0$, is not yet reviewed against the contract's code family.
\IRwhy An inconsistent or mistranscribed hypothesis silently breaks every conditional record.
Consistency cannot be proved in Lean.
\IRres Publish the policy: verbatim transcription over the exact code family, explicit constants, a
regime guard as a field, two independent reviewers, a catalogue of nearby false statements, and
revalidation on removal. Transcribe \lean{McaJohnson} as H-v1. The unconditional track uses UD,
which is proved in ArkLib \cite{BCIKS20}. T0 may open unconditional-only if this issue is late.
\IRmeta{cryptography reviewers, process}{P1}{1}{open}
\IRsrc{D1-G9, D1 §3.3; G-D3-3; R7 §8 item 3; R3 §6 item 3; R5 §7.1; DESIGN §5.}
\end{issue}

\begin{issue}{G-5}{Kernel, comparator and toolchain upgrades}
\IRprob Lean 4.34.0 fixed three routes to a proof of \lean{False}, and a nested-inductive
projection bug was accepted by an independent kernel too
\cite{lean4340,brossard2026leanproj}. There is no policy for security-relevant upgrades during a
season.
\IRwhy A known kernel bug on the pinned toolchain invalidates the trust story of every record.
\IRres
\begin{enumerate}[nosep,leftmargin=1.4em]
\item Monitor upstream advisories.
\item A security bump creates a new contract id and triggers automatic re-checking of all records
  (G-3); records that fail are withdrawn or ported.
\item Schedule bumps that are not security-relevant at season boundaries.
\item Keep a second kernel (nanoda) as defence in depth.
\end{enumerate}
\IRmeta{process, harness team}{P1}{1 (ongoing)}{open}
\IRsrc{\cref{subsec:thr-residual} item 3; R7 §8 item 6; R1 implications; \cref{subsec:arch-versioning}.}
\end{issue}

\begin{issue}{G-6}{Second-kernel policy (nanoda, \code{leanchecker})}
\IRprob CI replays the toy certificate in comparator's default kernel and in nanoda, and
\code{leanchecker} replays solution modules; debug options such as \lean{debug.skipKernelTC} are
caught only by these replays. No rule says which replays are mandatory for a verdict, which nanoda
version is pinned, what happens when the kernels disagree, or how the extra time (toy: 115\,s
default kernel, 175\,s with nanoda) is counted against the budget.
\IRwhy Kernel diversity is the only defence against single-kernel soundness bugs, and an undefined
disagreement procedure invites disputes over records.
\IRres Publish the rule: comparator's default kernel and nanoda (an independent implementation,
pinned by version and hash) are both mandatory; \code{leanchecker} replays in Lean's own kernel, so it
is an additional replay, not kernel diversity; a disagreement withholds the verdict
and is escalated as a kernel issue (G-5); replay time counts against the single-threaded budget.
\IRmeta{process, harness team}{P0}{0}{decision needed (both replays run in CI)}
\IRsrc{A4-F7; A6-N1; \cref{subsec:arch-ci}; \cref{tab:thr-catalogue} rows L3a, L7.}
\end{issue}

\begin{issue}{G-7}{Trusted-surface pins: unpinned binaries, host tools, elaborator and Lake}
\IRprob The trusted surface is wider than the Lean text and the \code{.olean} manifest: the Lean
elaborator and frontend for the challenge (the statement's meaning is the organiser-elaborated term),
Lake (which resolves the loaded \code{.olean}s; A6-G0 was a Lake behaviour), and the Lean compiler and
runtime (the gate, the axiom audit and the badge executables). The comparator, \code{lean4export} and
nanoda binaries are not covered by the manifest; nor are the host tools (Python~$\ge$3.11 for
\code{tomllib}, bubblewrap, \code{sha256sum}) and the harness shell scripts. Contract-used library
definitions (\lean{![…]}, \lean{List.ofFn}, \lean{Fin.cons}/\lean{append}/\lean{snoc}/\lean{induction},
\lean{Set}, \lean{BitVec}) need to be on the auditor's list (\cref{subsec:app-checklist}).
\IRwhy An unpinned checker binary or host tool is a silent change of the verdict function.
\IRres Pin the three checker binaries by hash under the contract id, in the manifest or the container
image (H-1); record the host-tool versions; extend the trusted-surface table accordingly.
\IRmeta{harness team, process}{P0}{0}{open}
\IRsrc{A7-M5; A6-G0; \cref{tab:arch-trusted}.}
\end{issue}

%-----------------------------------------------------------------------------------------------
\subsubsection{P2: T2, further tracks and research}

The nine P2 items block no season before their own; \cref{tab:iss-p2} gives each in one row.

\begingroup\small
\renewcommand{\arraystretch}{1.08}
\begin{longtable}{@{}l>{\raggedright\arraybackslash}p{0.34\textwidth}>{\raggedright\arraybackslash}p{0.36\textwidth}c>{\raggedright\arraybackslash}p{0.10\textwidth}@{}}
\caption{P2 issues (owners: contract team and cryptography reviewers unless noted).}\label{tab:iss-p2}\\
\toprule
ID & Issue & Proposed resolution & Phase & Sources \\
\midrule
\endfirsthead
\toprule
ID & Issue & Proposed resolution & Phase & Sources \\
\midrule
\endhead
\bottomrule
\endlastfoot
M-4 & Field-syscall board (T1-F) and precompiles: on plain RV64IM+Zbc, hand-tuned $E$ arithmetic dominates architecture & specify an \code{EMUL} system call ($E$ product on memory operands) priced from a binary-field zkVM, as a versioned machine extension with its own certificate variant and controls & 4 & A1-M5b, G-D3-11 \\
R-2 & Tree composition theorem \lean{tree\_secure} (L2) is unproved; it needs its own round count $k_{\mathrm{tree}}=Nk_e+(N-1)k_f+k_d$, the side condition $2k_{\mathrm{tree}}\le t$ and independent per-node challenges; until then T2 records certify one step & port leanerVM's \lean{KnowledgeAppend}, write the frame rule (\libc{3}), prove the theorem for the sequentialised protocol; harness enforces the declared depth (library team) & 3 & D1-G2, O14, A2-m4 \\
R-3 & All-runs cost bound for deciders: rejecting runs are bounded only by the fuel, $2^{36}$ steps per hart and $(\#\mathrm{spawns}+1)\cdot2^{36}$ with harts, a DoS exposure for deployed T2 deciders & planned, not in v0.3: \lean{AllRunsCostAtMost} as an optional displayed badge; from the T2 season (v0.4 or later) require $C^{\mathrm{all}}_{\mathrm{dec}}\le8\,C_{\mathrm{dec}}$ for T2 records, with library support for ranking arguments & 3 & G-D3-7, A6-F1 \\
R-4 & Bounded-depth accumulation \cite{BMNW24} is inadmissible in the depth-uniform interface & index $\Racc$ by level, $\mathrm{Fin}(d_{\max}+1)\to\mathrm{Set}$; each T2 record declares $d_{\max}$; tree theorem per level & 4 & D1-G13 \\
Q-1 & Efficient extraction (L2): at L1 the extractor's time is unconstrained, so L1 equals RBR soundness (\lean{SF.rbrKS\_of\_rbr}); the recursive bound consumes $t_E$ & \lean{ExtractorImage}, \lean{AgreesE} and a cycle bound $P_E$ of degree $\le3$, or a structural extractor; negative control: brute force fails the bound & 4 & D1-G5, A2-M2, A1-n2 \\
Q-2 & Recursive soundness after hash instantiation follows from no random-oracle theorem \cite{BCG24-RelROM} & keep \lean{RecursionHeuristic} an uninhabited schema, never a hypothesis of a scored theorem; track \cite{CGSY24,CCGOS23-AROM} (cryptography reviewers) & 3--4 & D1 §4.4 \\
Q-3 & Arithmetised recursion (T3): a trusted RV64IM+Zbc$\to$R1CS arithmetiser; the self-recursive fixed point needs $C\le1.7$--$5.6\times10^5$ cycles at $2^{24}$, and the best modelled design is 9--29 times above it & research: verified IR compiler to RV64IM (5--15k lines, A3-S7) or a field-native machine (T1-F) as stepping stones & 5 & D2-G13, A3-S7 \\
Q-4 & ROM argument track (L3): BCS engineering (caps, arity, pruning, grinding; 1.5--2$\times$ of real hash cost) is invisible in the IOP model & open the \code{HASH} call (number 2 is reserved) and Zbb, with sig.golf-style grinding budgets; needs a generic BCS theorem \cite{BCS16,CY24-Book} & 4 & D2-G6, A1-m3 \\
Q-5 & Further tracks: Boolean $\mathrm{GF}(2)$ R1CS with packed witnesses; certified attacks on parameter sets & the same contract with $\F:=\mathrm{GF}(2)$, a packing clause and its own baseline; a Lean lower-bound statement form for attacks & 4 & G-D3-9, D3 §1.5 \\
\end{longtable}
\endgroup

%===============================================================================================
\subsection{Findings resolved in DESIGN v0.2 and v0.3}\label{subsec:iss-resolved}

\Cref{tab:iss-resolved} maps each finding that the specification in force answers to the place
where the answer is given: a row of the attack catalogue (\cref{tab:thr-catalogue}) or a section of
this report. V1--V36 were resolved in DESIGN v0.2; V37--V55 in the compiled v0.3 contract, the
round-3 and round-4 harness and the certified baseline. An open remainder is named by its issue in parentheses.
The twelve points of the brief's review are mapped in \cref{tab:brief-review}.

\begingroup\small
\renewcommand{\arraystretch}{1.02}
\begin{longtable}{@{}l>{\raggedright\arraybackslash}p{0.37\textwidth}>{\raggedright\arraybackslash}p{0.19\textwidth}>{\raggedright\arraybackslash}p{0.20\textwidth}l@{}}
\caption{Resolved findings and where they are answered.}\label{tab:iss-resolved}\\
\toprule
ID & Finding & Source & Answered by & DESIGN\,§\\
\midrule
\endfirsthead
\toprule
ID & Finding & Source & Answered by & DESIGN\,§\\
\midrule
\endhead
\bottomrule
\endlastfoot
V1 & Typed promise: submitter message types give free proximity & A1-C1, A2-C1 & S1 & 3, 4 \\
V2 & Submitter challenge types; unpriced FS sampling; tape encoding & A1-M1, A2-C2 & S2 (constants: M-3) & 3, 5, 7 \\
V3 & No round cap; GKR nearly banned & A2-C3, A1-M6 & S3 & 3 \\
V4 & Knowledge headline cites the wrong theorem; side conditions & A2-M1, A2-m1, D1-G4 & \cref{subsec:sec-cap} & 5 \\
V5 & L1 \enquote{knowledge} equals RBR soundness & A2-M2, A1-n2, R5 §7.3 & S7 (L2: Q-1) & 5 \\
V6 & Decorative accumulator caps and interface & A2-M3, A1-m4 & S4--S6 & 6 \\
V7 & Indexer may be noncomputable & A2-M4, A3-m2 & S8 (scope: C-7) & 3, 8 \\
V8 & Record gate bypass through VK/IDX & A1-M3 & H1 (build: H-5) & 8 \\
V9 & $n_x\le2^{16}$ public-input floor & A1-M2 & C6 & 2 \\
V10 & Prover cap a corner solution; prover cost unobservable & A1-M4, D1-G6 & C5 & 7 \\
V11 & One scalar cost inverts real recursion & A1-M5a & C14, \cref{subsec:cost-vector} & 7 \\
V12 & Three incompatible machines & A3-C1, D1-G3 & \cref{subsec:proto-machine} (M-1) & 4 \\
V13 & Cost and termination over adversarial indices & A3-C2, A3-n2, A3-S1/S2, U2 & S13, C2 & 4, 7 \\
V14 & 4\,GiB memory too small at $2^{24}$ & A3-M1 & $2^{36}$-byte memory & 4 \\
V15 & Budgets stated in lines & A3-M4, R2 §6.3 & \cref{subsec:arch-budgets} & 11 \\
V16 & T0 has no query semantics & A3-M5, A3-S6 & \lean{QProg} (C-4) & 11 \\
V17 & Toolchain mismatch & A3-M6, R5, R6, PLAN §1.11 & \cref{subsec:arch-toolchain} (G-2) & 13 \\
V18 & Prime-field limits at $2^{24}$; error cap vs extension & U1, D2-G7/G8/G15, A1-m9 & \cref{subsec:proto-fields}, S12 & 2, 5 \\
V19 & Binary fields need carry-less multiply & R8 §5, D1-G8, D2-G4, A1-m7 & Zbc in the ISA, priced per profile (M-1) & 4, 15.1 \\
V20 & QUERY granularity and pricing & D1-G7, D2 §9.2 & C8, C25 & 4, 7 \\
V21 & Free Merkle caps & A1-m3 & C7 (Q-4) & 7 \\
V22 & $/64$ bans linear deciders & A1-m5, A5-M8 & C10 & 15.3 \\
V23 & $\lambda=1/8$ is dead weight & A1-m1 & \cref{subsec:cost-zk} & 7 \\
V24 & Ambiguities: $N$, IDX in $\sum\ell$, tape cap, leaf bytes & A1-m10 & \cref{subsec:cost-zk} & 7 \\
V25 & Memory first-touch cost & D2-G3, A1-m6 & C9 & 15.1 \\
V26 & \enquote{Performance $\div$ size}; median over runs & PLAN §1.2, §1.9 & \cref{subsec:cost-candidates} & 1 \\
V27 & Overfitting to the benchmark index & PLAN §1.6, D1-G12, A1-n1 & C1 (B-1) & 9 \\
V28 & Proximity conjectures & PLAN §1.8, R3, R7 & S11 (G-4) & 5 \\
V29 & FS statement binding; sampler bias & R7, A2-m5, D1 §3.1, R8 & S14 (H-3) & 5 \\
V30 & ArkLib \code{sorry} taint & A2-m2, R4 §9 & L11 (L-8) & 5, 13 \\
V31 & Certificates for every size below the cap & A3-S3 & class shapes and padding & 2 \\
V32 & Irreducibility certificates may break replay & G-D3-2, R3, R8 & L8 & 2 \\
V33 & Fuel vs $2^{24}$ demonstration time & G-D3-12, A3 §6.3 & \cref{subsec:cost-vector} (H-4) & 4 \\
V34 & Batcher property is completeness only & PLAN §1.7, R5 §7.2 & \cref{sec:recursion} & 6 \\
V35 & Quantum adversaries & D1-G14 & out of scope & 5 \\
V36 & KoalaBear \lean{norm\_num} probe failed & A3-n3 & obsolete (binary field) & 2 \\
\midrule
V37 & Dense \lean{Index}: honest indexer infeasible & I-BAS-1 & sparse \lean{Index} (\cref{subsec:base-findings}) & 15 \\
V38 & Padding must reflect satisfiability (former C-2) & O13, A3-S3 & \lean{pad\_sat\_iff} & 2 \\
V39 & Untyped accumulator oracles & A4-F1 & S16 & 15.2 \\
V40 & Image and score inputs unbound & A4-F2 & S17 & 15.2 \\
V41 & Text-scan hygiene gate bypassable & A4-F3 & L3a & 15.4 \\
V42 & Score not a Lean function & A4-F4 & S18 & 15.2 \\
V43 & Permissive claim schema & A4-F5 & H9 & 15.2 \\
V44 & FS absorption of $\vk\,\|\,x$ unpriced & A4-F6 & C24 & 15.2 \\
V45 & Weak CI: fake sandbox, one kernel & A4-F7 & H10 (H-7, G-6) & 15.4 \\
V46 & Build order lets comparator accept a false statement & A6-G0 & L12 & 15.4 \\
V47 & Leaf bytes at v0.2's hash rate & A6-P1 & C25 (latency: M-6) & 15.4 \\
V48 & \lean{@[init]} passes the gate & A6-G1 & L13 & 15.4 \\
V49 & Gate trusts rewritable \code{.olean}s by location & A6-G2 & L14 & 15.4 \\
V50 & Renderer reads any file the claim names & A6-R1 & H11 & 15.4 \\
V51 & Baseline end to end on the contract to be frozen (former L-3) & O9, A3-C1 & \cref{sec:baseline}: v0.3 T1 and \lean{Latency1} certified, HARNESS: ACCEPT & 10, 15 \\
V52 & Compile the accumulation interface (former R-1) & A2-M3, O8 & S6; \code{Accumulation.lean}, controls C5, A2-C & 6, 15.2 \\
V53 & Comparator end to end on contract plus baseline (former H-2) & A3-M3 & comparator stage 243.6\,s with both kernels (\cref{tab:base-time}) & 13 \\
V54 & Gate rule~4 rejected every protocol whose declared error uses \lean{ℝ≥0} operations & A7-M2, A6-N3 & rule~4 checks only the badge-path fields (\code{probes/a7/m2}) & 3 \\
V55 & Gate and badge ran on a second, uncompared build & A7-M3 & L15: gate-tree build byte-identical to the comparator-verified one (\code{bindproj}) & 15.4 \\
\end{longtable}
\endgroup

%===============================================================================================
\subsection{Roadmap}\label{subsec:iss-roadmap}

The roadmap orders the open issues into six phases. Each phase ends with explicit go/no-go
criteria. Effort figures are A3's line estimates, converted at A3's implied rate of about 2k
library lines per organiser agent-week, which is calibrated on \libc{0}--\libc{2}. They are
indicative. Calendar months in \cref{fig:iss-gantt} assume parallel work streams, and the
dependencies matter more than the dates.

\begingroup\small
\renewcommand{\arraystretch}{1.12}
\begin{longtable}{@{}>{\raggedright\arraybackslash}p{0.13\textwidth}>{\raggedright\arraybackslash}p{0.30\textwidth}>{\raggedright\arraybackslash}p{0.16\textwidth}>{\raggedright\arraybackslash}p{0.33\textwidth}@{}}
\caption{Phased plan. Issue identifiers refer to \cref{subsec:iss-open}.}\label{tab:iss-roadmap}\\
\toprule
Phase & Scope (issues) & Effort [indicative] & Go/no-go criteria for exit \\
\midrule
\endfirsthead
\toprule
Phase & Scope (issues) & Effort [indicative] & Go/no-go criteria for exit \\
\midrule
\endhead
\bottomrule
\endlastfoot
0. Contract freeze &
v0.3 contract with the round-3 fixes and its controls (C-1, C-3, C-4); Zbc tests (M-1); class decision (M-2); \libc{0}, \libc{1} (L-1, L-2); pipeline, intake, container and production sandbox (H-1, H-7, H-9); trusted-surface pins (G-7); second-kernel policy (G-6); final review round and bounty (G-1); release and versioning (G-2, G-3) &
\libc{0}--\libc{1}: 4.5--7.5k lines, 3--5 agent-weeks; contract, harness stages 0--2 and the baseline done (v0.3, round~3); review 1--2 reviewer-weeks; intake about 1k lines (port) &
(i)~every negative control fails and every positive control holds (\cref{tab:sec-controls});
(ii)~\emph{met:} the baseline is certified on the v0.3 machine, and the organiser's verdict (one harness run, both kernels) takes $\le10$\,min warm and $\le60$\,min cold; measured 6.1\,min per run, 13.6\,min for the submitter's full check of two records (\cref{tab:base-time});
(iii)~\emph{met:} comparator export and replay of contract plus baseline take $\le10$\,min (comparator stage 243.6\,s with both kernels), and comparator rejects an altered contract definition and \lean{native\_decide};
(iv)~trusted contract $\le850$ code lines excluding vendored files (v0.3: 796); vendor check passes; \\
(Phase 0, cont.) & & &
(v)~no open critical or high finding from the review rounds (A4, A6 and the round on the release candidate);
(vi)~a clean-room rebuild from the container reproduces the contract id and the \code{.olean} manifest;
(vii)~the harness enforces the order of Decision~A-3 (comparator first in a fresh tree, gate in a second fresh tree, badge as its own process) under the production sandbox, with the organiser package read-only and the manifest checked before and after every stage; the \code{cmpproj} and \code{gateproj} regressions are rejected;
(viii)~the contract is frozen at v0.3 plus the round-3 fixes; the v0.4 items (C-5, M-5--M-8) are scheduled, not merged. \\
\midrule
1. T0 season (\lean{QProg}) &
T0 board open; \libc{2}, \libc{3} (L-4, L-5); bridge (L-8); executable $K$/$E$ (L-9); hart combinators (L-10); union bound (C-6); constants (M-3); v0.4 contract: fuel rule, spawn pricing, latency leaf, timer hazards, leaf padding (C-5, M-5--M-8); feasibility rules (C-7); FS specification (H-3); emulator (H-4); badge path (H-5); image admission (H-6); generators (B-1, B-2); hypothesis file (G-4); upgrade policy (G-5); start \libc{4} (L-6) &
\libc{2}--\libc{3}: 5--11k lines, 3--6 agent-weeks; \libc{6} rest 3.5--5.5k lines; reviewers for G-4, M-3 &
(i)~\libc{0}--\libc{2} shipped, which is the opening condition of T1 in DESIGN §11;
(ii)~emulator differential tests show zero mismatches, and the $2^{24}$ baseline run takes $\le15$\,min;
(iii)~measured $\le$ claimed for the baseline on every benchmark family and class;
(iv)~generators published with hashes and cross-checked against the baseline;
(v)~a badge-path dry run on the baseline passes;
(vi)~FS constants measured and the scorer version frozen for the T1 season;
(vii)~v0.4 released under a new contract id with C-5 and M-5--M-8 decided, records revalidated (G-3), and the round-3 regressions still rejected. \\
\midrule
2. T1 launch &
T1 board open; \libc{4} interim then full (L-6); \libc{5} reference submission (L-7) &
\libc{4}: interim 5--10k lines (3--5 agent-weeks), full 15--40k (8--20); \libc{5}: 10--20k (5--10) &
(i)~\libc{5} certified within the T1 budget (3\,h, 8 CPUs, 48\,GiB);
(ii)~it passes the badge path (prover $\le1$\,h at $2^{20}$, $\le8$\,h at $2^{24}$);
(iii)~its $\mathrm{Score}_{zk}$ is below the certified baseline's $1.624\times10^9$ (\cref{tab:base-score});
(iv)~the \libc{4} PCS has a sorry-free RBR-KS theorem (interim flagged until the full one exists). \\
\midrule
3. T2 accumulation &
tree theorem (R-2); decider all-runs clause (R-3); \lean{RecursionHeuristic} text (Q-2) &
tree theorem on top of \libc{3}'s append and frame rule; 1--3k lines [judgement] &
(i)~\lean{tree\_secure} proved with its own $k_{\mathrm{tree}}$, independent per-node challenges and the $2k_{\mathrm{tree}}\le t$ condition;
(ii)~\lean{fullRecAcc} and one non-trivial accumulation certified;
(iii)~T2 labels switch from \enquote{one-step} to \enquote{tree} on re-check. \\
\midrule
4. Further tracks &
L2 efficient extraction (Q-1); ROM track (Q-4); Boolean track and attack track (Q-5); bounded depth (R-4); T1-F and precompile board (M-4) &
each track: contract extension, controls, baseline; ROM needs a generic BCS theorem [large] &
per track: its own positive and negative controls pass, including that a brute-force extractor fails the L2 cycle bound; a baseline is certified; the extension gets a new contract id, with records revalidated. \\
\midrule
5. T3 arithmetised recursion &
arithmetiser or verified IR compiler; fixed-point study (Q-3) &
research; compiler 5--15k lines (A3-S7); arithmetiser of SP1/OpenVM scale &
an arithmetiser correctness theorem against the contract machine, and a modelled fixed point inside the $2^{24}$ class. \\
\end{longtable}
\endgroup

\begin{figure}[t]
\centering
\begin{tikzpicture}[x=0.56cm,y=0.42cm,
  bar/.style={draw=black!55,rounded corners=1pt,minimum height=0.30cm,inner sep=0pt},
  p0/.style={bar,fill=red!30}, p1/.style={bar,fill=orange!35}, p2/.style={bar,fill=blue!25},
  p3/.style={bar,fill=green!30}, p4/.style={bar,fill=violet!25}, p5/.style={bar,fill=black!12},
  lab/.style={anchor=east,font=\scriptsize,inner sep=1pt},
  dep/.style={-{Stealth[length=1.6mm]},draw=black!70,thin},
  ms/.style={draw=black,fill=black,diamond,inner sep=1.2pt}]
% axis
\foreach \m in {0,2,...,18} {
  \draw[black!15] (\m,0.6) -- (\m,-19.6);
  \node[font=\scriptsize,anchor=south] at (\m,0.6) {\m};
}
\node[font=\scriptsize,anchor=south east] at (-0.15,0.6) {month};
% helper: \bar{row}{start}{end}{style}{label}
\newcommand{\gbar}[5]{%
  \draw[#4] (#2,-#1+0.17) rectangle (#3,-#1-0.17);
  \node[lab] at (-0.15,-#1) {#5};}
\gbar{0}{0}{2.5}{p0}{C-1/C-4 contract v0.3 + controls}
\gbar{1}{0}{1.5}{p0}{\libc{0} machine facts (L-1)}
\gbar{2}{0.5}{2.5}{p0}{\libc{1} executor (L-2)}
\gbar{3}{0.5}{3}{p0}{intake, preflight, pins (H-9, G-7)}
\gbar{4}{0}{3}{p0}{harness, container, sandbox (H-1, H-7)}
\gbar{5}{2}{3}{p0}{review, bounty, freeze (G-1, G-2, G-6)}
\gbar{6}{3}{7}{p1}{T0 season (\lean{QProg})}
\gbar{7}{2}{5}{p1}{\libc{2} field routines (L-4)}
\gbar{8}{3}{7}{p1}{\libc{3} IOP components (L-5)}
\gbar{9}{3}{6}{p1}{emulator, badge path, generators (H, B)}
\gbar{10}{3}{6}{p1}{v0.4 (C-5, M-5--M-8)}
\gbar{11}{3}{7}{p2}{\libc{4} interim tensor PCS (L-6)}
\gbar{12}{4}{12}{p2}{\libc{4} binary RS PCS, critical path (L-6)}
\gbar{13}{6}{12}{p2}{\libc{5} reference submission (L-7)}
\gbar{14}{6}{18}{p2}{T1 season}
\gbar{15}{9}{13}{p3}{tree theorem (R-2)}
\gbar{16}{12}{18}{p3}{T2 season}
\gbar{17}{11}{18}{p4}{L2 extraction, ROM, Boolean, T1-F (Q, M-4)}
\draw[p5,dashed] (14,-18+0.17) rectangle (18,-18-0.17);
\node[lab] at (-0.15,-18) {T3 research (Q-3)};
% milestones
\node[ms] (fz) at (3,-19) {}; \node[font=\scriptsize,anchor=north] at (3,-19.15) {freeze, T0 opens};
\node[ms] (t1) at (6,-19) {}; \node[font=\scriptsize,anchor=north] at (6.4,-19.15) {T1 opens};
\node[ms] (t2) at (12,-19) {}; \node[font=\scriptsize,anchor=north] at (12,-19.15) {T2 opens};
% dependencies
\draw[dep] (1.5,-1) -- (1.5,-1.8);
\draw[dep] (2.5,-2) .. controls (2.9,-3) .. (3.05,-4.9);
\draw[dep] (5,-7) .. controls (5.6,-8) and (5.8,-12) .. (6,-13.85);
\draw[dep] (7,-11) -- (7,-12.85);
\draw[dep] (12,-12) -- (12,-15.85);
\draw[dep] (13,-15) -- (13,-15.85);
\draw[dep] (7,-8) .. controls (8.5,-9) and (9,-13) .. (9.1,-14.85);
% phase legend
\begin{scope}[shift={(0,-21.3)}]
\foreach \s/\t/\x in {p0/Phase 0/0, p1/Phase 1/3, p2/Phase 2/6, p3/Phase 3/9, p4/Phase 4/12, p5/Phase 5/15} {
  \draw[\s] (\x,0.15) rectangle (\x+0.8,-0.15);
  \node[font=\scriptsize,anchor=west] at (\x+0.9,0) {\t};
}
\end{scope}
\end{tikzpicture}
\caption{Indicative schedule with the main dependencies (arrows): \libc{0}$\to$\libc{1}$\to$freeze; \libc{2}$\to$T1 opens; interim \libc{4}$\to$\libc{5}; full \libc{4}$\to$T2;
\libc{3}$\to$tree theorem$\to$T2. Months count from the start of Phase~0.}
\label{fig:iss-gantt}
\end{figure}

\paragraph{Reading the plan.} Phase~0 is the only phase whose output is trusted; after the freeze
only security fixes change the contract id (G-2). Its exit criteria include the harness order and the
manifest (vii), because round~3 showed that a correct contract can still certify a false statement
under a wrong build order. The v0.4 items (\cref{tab:iss-v04}) land before T1 with a new
contract id. Phase~1 overlaps the T0 season with the library work T1 needs (A3-S6). Until \libc{4} and
\libc{5} exist, T1 is expected to show only trivial and hybrid entries, so \libc{4}, the critical path,
starts in Phase~1 with the interim tensor-code PCS as a hedge. Phases 4 and 5 enter as versioned
contract extensions with their own controls and baselines.

\FloatBarrier
\begin{takeaway}[Launch-blocking work in one paragraph]
The 15 P0 issues are C-1, C-3, C-4, M-1, M-2, L-1, L-2, H-1, H-7, H-9, G-1, G-2, G-3, G-6 and G-7.
The certified baseline already shows the certificate inhabited and checkable in about 6\,min per
harness run. Before anything opens:
\begin{itemize}[nosep]
\item freeze the v0.3 contract with the round-3 fixes and all controls, including the T0
  certificate; publish the compiled-code policy that the gate already enforces;
\item finish the Zbc differential tests, decide the class set and ship \libc{0}--\libc{1};
\item build the intake (preflight, source policy) and run the harness in the order of Decision~A-3
  under the production sandbox, with the \code{.olean} manifest checked before and after every stage
  and the checker binaries pinned; fix the second-kernel policy;
\item have the release candidate, pipeline included, reviewed once more and offer a bounty;
\item freeze it with versioning and revalidation in place.
\end{itemize}
Everything else, including the v0.4 cost and timer fixes and the critical-path PCS, can follow the T0
launch.
\end{takeaway}
\clearpage
